\documentclass[11pt]{article}
\usepackage[utf8]{inputenc}
\usepackage{amsmath}
\usepackage{geometry}
\geometry{letterpaper, margin=1in}

\title{\textbf{Editorial: Problem F - Rewrite The Stars}}
\author{Setter: Hudson}
\date{}

\begin{document}
\maketitle

\section*{Problem Overview}
We are tasked with maintaining the \textbf{Maximum Weight Independent Set (MWIS)} of a dynamic forest under the following operations:
\begin{enumerate}
    \item \textbf{Link($u, v$):} Add an edge between two distinct trees.
    \item \textbf{Cut($u, v$):} Remove an existing edge.
    \item \textbf{Update($u, w$):} Update the weight of vertex $u$.
    \item \textbf{Query($u$):} Return the MWIS of the connected component containing $u$.
\end{enumerate}

\section*{Difficulty Analysis}
Rating: \textbf{3500+ (Legendary Grandmaster)}.
This problem represents the pinnacle of dynamic tree data structures and is likely the hardest problem in the contest.
\begin{itemize}
	\item Standard \textbf{Dynamic DP} (via Heavy-Light Decomposition) can handle weight updates but cannot handle structural changes (Link/Cut).
	\item Standard \textbf{Link-Cut Tree} (LCT) handle structural changes but are designed for path aggregates. They cannot easily maintain subtree aggregates (like MWIS) because the "dashed edges" in an LCT are not dynamically aggregated in a way that supports non-invertible DP transitions efficiently without complex auxiliary structures.
\end{itemize}
To solve this, one must implement a \textbf{Self-Adjusting Top Tree} (specifically the Compress-Rake Cluster variant). This data structure generalizes LCTs to handle subtree aggregates dynamically by treating "Rake" operations (merging a subtree onto a path) as first-class citizens alongside "Compress" operations (merging paths).

\section*{Solution Approach}
The core idea is to decompose the tree into \textbf{Clusters}. A cluster corresponds to a connected subgraph with at most two \textbf{Boundary Nodes} interacting with the rest of the tree.

For MWIS, we define the state of a cluster $(u, v)$ (where $u, v$ are boundaries) using a $2 \times 2$ matrix $M$:
$$ M[a][b] = \text{Max weight IS of cluster s.t. } u \text{ has state } a, v \text{ has state } b $$
State $0 \implies$ Excluded, State $1 \implies$ Included.

\subsection*{Top Tree Operations}
The Top Tree maintains the structure using two merge types:
\begin{enumerate}
    \item \textbf{Compress (Path Merge):} Merges two clusters $(u, k)$ and $(k, v)$ into $(u, v)$.
    The transition resembles matrix multiplication in the $(\max, +)$ semiring.
    $$ M_{new}[a][b] = \max(L[a][0] + R[0][b], L[a][1] + R[1][b] - W_k) $$
    \textit{(Note: $W_k$ is subtracted to prevent double-counting the shared boundary node $k$.)}
    
    \item \textbf{Rake (Subtree Merge):} Merges a cluster $(k, k)$ (effectively a pending subtree) onto a path cluster at node $k$.
    This updates the internal weight/state of node $k$ within the path cluster to reflect the optimal choice from its subtree.
\end{enumerate}

\subsection*{Implementation Complexity}
This solution requires implementing a full Top Tree with:
\begin{itemize}
    \item Lazy propagation (for path reversals).
    \item Ternary Splay operations (handling Left Compress, Right Compress, and Middle Rake children).
    \item `Access`, `Link`, and `Cut` operations that explicitly manage Rake edges.
\end{itemize}
A correct implementation typically exceeds 200 lines of C++. Due to the complexity, the time limit is set generously to 8.0s.

\end{document}
